\documentclass[11pt]{article}
\usepackage[a4paper,margin=1in]{geometry}
\usepackage{subfiles}
\usepackage{amsmath,amssymb,bm}
\usepackage{hyperref}
\usepackage{booktabs}
\usepackage{array}
\usepackage{mathtools}

\newcommand{\rs}{r_{\mathrm{s}}}

\title{Appendix 6 (to main theory): Structure of the 2PN Shapiro Delay in the Bi--Conformal ISPG Metric}
\author{}
\date{}

\begin{document}
\let\phi\varphi
\maketitle

\section*{Purpose and scope}
This appendix fixes the \emph{structural} form of the
second post-Newtonian (2PN) Shapiro time delay in the
bi--conformal ISPG geometry.  It proves one finite,
convention-independent sub-piece---the \emph{refractive}
contribution---as an exact analytic statement, and it
sets up (but does not close) the two remaining pieces
that together compose the full 2PN coefficient:

\begin{enumerate}
\item a \textbf{refractive (straight-line) contribution} from
the explicit $O(\rs^2)$ term in the bi-conformal refractive
index, which is derived exactly and verified numerically;
\item a \textbf{bent-ray refractive contribution} from the
$O(\rs)$ light-bending correction inserted into the $O(\rs)$
time integrand, whose \emph{finite} coefficient depends on
consistent choices of (a) asymptotic impact parameter vs.\
geometric closest approach and (b) finite endpoint
regularization $r_1,r_2$;
\item an \textbf{excess path-length contribution} from the
Euclidean length difference between the bent ray and the
straight-line reference, which is infra-red sensitive in the
same sense as (ii).
\end{enumerate}

\paragraph{Epistemic status.}
Piece~(i) is established here analytically and has been
cross-checked numerically
(\texttt{MAIN/Python/Shapiro/check\_shapiro\_born.py}).  Pieces~(ii)
and~(iii) are flagged as \emph{in progress}: their individual
finite coefficients are not asserted, and the full 2PN
coefficient $\Delta t_{2\mathrm{PN}}^{\mathrm{ISPG}}$ is left
as future work.  Although the individual pieces remain open,
their dependence is entirely on the shared $O(r_g)$ trajectory
of the two metrics, so they cancel in the ISPG--GR
\emph{difference}; the difference of the full 2PN coefficient
is therefore finite, IR-insensitive, and computed exactly in
Sec.~\ref{sec:diff}, with the closed form
\[
\Delta t_{2{\rm PN}}^{\rm ISPG}-\Delta t_{2{\rm PN}}^{\rm GR}
=\frac{r_g^2}{cb}\cdot\frac{\pi}{4}.
\]
The main theory Shapiro section~\cite{Lotishvili2026} carries the structural statement
that the bi-conformal exponential refractive index
$n=e^{-\phi}$ produces a 2PN Shapiro signature distinct from
GR's rational-polynomial form; the closed-form ISPG--GR
discriminator above is the operational expression of that
distinction.

\section{Conventions and the static ISPG exterior}
We use the static, spherically symmetric ISPG line element (the main theory)
\begin{equation}
ds^2=-e^{\phi(r)}c^2dt^2+e^{-\phi(r)}\left(dr^2+r^2d\Omega^2\right),
\label{eq:metric_app6}
\end{equation}
and we define the weak-field length scale by
\begin{equation}
\phi(r)\simeq -\frac{2GM}{c^2 r}\equiv -\frac{2r_g}{r},
\qquad r_g\equiv \frac{GM}{c^2}.
\label{eq:phi}
\end{equation}
(With this convention the Newtonian potential $\Phi_N=-GM/r$ corresponds to $g_{tt}\simeq -(1+2\Phi_N/c^2)=-(1+\phi)$.)

\section{Hamilton--Jacobi solution for null geodesics}
For equatorial motion ($\theta=\pi/2$), take the separated action
\begin{equation}
S=-Et+L\varphi+S_r(r),
\end{equation}
and impose the HJ equation $g^{\mu\nu}\partial_\mu S\,\partial_\nu S=0$ for null rays. With \eqref{eq:metric_app6} this yields
\begin{equation}
p_r^2\equiv \left(\frac{dS_r}{dr}\right)^2=\frac{E^2}{c^2}e^{-2\phi(r)}-\frac{L^2}{r^2}.
\label{eq:pr2}
\end{equation}
The asymptotic impact parameter is the invariant ratio
\begin{equation}
b\equiv \frac{Lc}{E},
\label{eq:bdef}
\end{equation}
defined in the Minkowski limit $\phi\to 0$.

Throughout this appendix we take~\eqref{eq:bdef} as the operational
definition of the impact parameter; the distinct quantity
$r_{\min}$ (geometric closest approach of the bent trajectory)
differs from $b$ at $O(r_g)$, and the distinction is essential
to the finite-coefficient determination in Sec.~\ref{sec:open}.

The zeroth-order trajectory is the straight line
\begin{equation}
r_0(z)=\sqrt{b^2+z^2}
\label{eq:r0_app6}
\end{equation}
along the asymptotic direction $z$.

\section{Time-of-flight integral and its 2PN decomposition}
From $ds^2=0$ in \eqref{eq:metric_app6},
\begin{equation}
0=-e^{\phi}c^2dt^2+e^{-\phi}dl^2
\quad\Rightarrow\quad
dt=\frac{e^{-\phi(r)}}{c}\,dl,
\label{eq:dt_app6}
\end{equation}
where $dl$ is the Euclidean line element in the conformal
spatial metric (the bracket in~\eqref{eq:metric_app6}).
The observable Shapiro delay is the excess over the
flat-space straight-line travel time between the same
asymptotic endpoints:
\begin{equation}
\Delta t \equiv \frac{1}{c}\int_{\rm ray}\left(e^{-\phi(r)}-1\right)dl
  + \frac{L_{\rm ray}-L_0}{c}\,,
\label{eq:dtdef}
\end{equation}
where $L_{\rm ray}$ is the Euclidean length of the bent
spatial path and $L_0$ the straight-line length.
With \eqref{eq:phi}, the exponential expands as
\begin{equation}
e^{-\phi}=e^{2r_g/r}=1+\frac{2r_g}{r}+\frac{2r_g^2}{r^2}+O(r_g^3).
\label{eq:exp}
\end{equation}
At 2PN order ($O(r_g^2)$) three contributions appear:
\begin{align}
\Delta t_{2{\rm PN}} &= \Delta t_{2{\rm PN}}^{\rm (refr)}+\Delta t_{2{\rm PN}}^{\rm (bent)}+\Delta t_{2{\rm PN}}^{\rm (path)}, \label{eq:split_app6}\\
\Delta t_{2{\rm PN}}^{\rm (refr)}&=\frac{1}{c}\int \frac{2r_g^2}{r_0(z)^2}\,dl_0, \label{eq:refr2pn}\\
\Delta t_{2{\rm PN}}^{\rm (bent)}&=\frac{1}{c}\int \frac{2r_g}{r}\,dl-\frac{1}{c}\int \frac{2r_g}{r_0}\,dl_0\Bigg|_{O(r_g^2)},
\label{eq:bent2pn}\\
\Delta t_{2{\rm PN}}^{\rm (path)}&=\frac{L_{\rm ray}-L_0}{c}\Bigg|_{O(r_g^2)}.
\label{eq:path2pn}
\end{align}
Here $dl_0\simeq dz$ and $r_0(z)$ is the straight path \eqref{eq:r0_app6}.

\section{Refractive-source term (derived)}
Using $dl_0=dz$ and $r_0^2=b^2+z^2$,
\begin{equation}
\Delta t_{2{\rm PN}}^{\rm (refr)}
=\frac{2r_g^2}{c}\int_{-\infty}^{+\infty}\frac{dz}{b^2+z^2}
=\frac{2r_g^2}{c}\cdot \frac{\pi}{b}
=\frac{r_g^2}{cb}\,(2\pi).
\label{eq:refr2pi}
\end{equation}
Thus the quadratic term of the exponential expansion of the
bi-conformal refractive index contributes \emph{exactly}
$2\pi$ in the coefficient multiplying $r_g^2/(cb)$.  This
piece is convention-independent (it involves only the straight
reference line and the exactly-computable integral
$\int dz/(b^2+z^2)=\pi/b$).

\paragraph{Numerical verification.}
The integral~\eqref{eq:refr2pi} has been cross-checked by
direct numerical evaluation of
$\int_{-Z}^{+Z}(e^{2r_g/r_0}-1)\,dz$ at several $r_g/b$ ratios
and finite cutoffs~$Z$.  Extracting the $O(r_g^2)$ coefficient
by polynomial regression in $r_g$ (with $b=1$, $Z\to 2000\,b$)
gives $B\to 6.281\ldots\to 2\pi$.  The driver script is
\texttt{MAIN/Python/Shapiro/check\_shapiro\_born.py}.

\paragraph{Structural comparison with GR.}
The GR Schwarzschild isotropic metric yields a refractive
index that is a rational function of $r_g/r$ rather than an
exponential.  The analogous $O(r_g^2)$ straight-line integrand
therefore differs from~\eqref{eq:refr2pi}, and the refractive
sub-piece in GR has a different numerical coefficient.  This
structural difference (exponential vs.\ rational) is the
basic origin of any 2PN discrepancy between ISPG and GR and
survives independently of the two remaining pieces below.

\section{Bent-ray refractive and path-length contributions (in progress)}\label{sec:open}

The two remaining 2PN contributions involve the
$O(r_g)$ trajectory correction and are currently \emph{open}
in this appendix.  We record here only the mathematical
framework and the issues that must be resolved for a rigorous
finite-coefficient statement.

\subsection{$O(r_g)$ trajectory correction: convention issue}

From~\eqref{eq:pr2} and~\eqref{eq:phi}, expanding $p_r$ to first
order in $r_g$ and writing the trajectory equation in the
$u\equiv 1/r$ form produces the standard weak-field bending
equation
\begin{equation}
\frac{d^2u}{d\varphi^2}+u=\frac{2r_g}{b^2}+O(r_g^2),
\label{eq:ubending}
\end{equation}
whose solution with asymptotic straight-line boundary
conditions is
\begin{equation}
u(\varphi)=\frac{\cos\varphi}{b}+\frac{2r_g}{b^2}+O(r_g^2).
\label{eq:usol}
\end{equation}
This yields the familiar $O(r_g)$ deflection
$\delta\theta=4r_g/b$ and, at closest approach $\varphi=0$, the
geometric offset $r_{\min}=b-2r_g+O(r_g^2)$ \emph{relative to}
the HJ impact parameter~\eqref{eq:bdef}.

A straight-line-parameterized expression for the trajectory
correction $\delta r(z)$ depends on which of two distinct
quantities one identifies with $b$:
\begin{itemize}
\item If $b$ denotes the HJ asymptotic impact parameter
$Lc/E$~\eqref{eq:bdef}, then at $z=0$ the correction is
$\delta r(0)=-2r_g$ (the ray dips inside the flat-space straight
line by $2r_g$ at closest approach).
\item If $b$ denotes the geometric closest approach (so that
$r(0)=b$ by construction), then by definition $\delta r(0)=0$,
and the asymptotic impact parameter is shifted to $b+2r_g$.
\end{itemize}
The two conventions describe the \emph{same physical
trajectory} but differ in where the labelled parameter~$b$
sits, and they therefore produce different \emph{intermediate}
expressions for the contributions below.  A consistent 2PN
derivation must fix one convention from the outset and apply
it uniformly across~\eqref{eq:refr2pn}--\eqref{eq:path2pn}.

\subsection{Bent-ray refractive and path-length integrals}

Inserting $r=r_0+\delta r$ into the $O(r_g)$ integrand
$(2r_g/r)$ and keeping only the $O(r_g^2)$ contribution
produces
\begin{equation}
\Delta t_{2{\rm PN}}^{\rm (bent)}
\;=\;-\frac{2r_g}{c}\int \frac{\delta r(z)}{r_0(z)^2}\,dz
\;+\;O(r_g^3),
\label{eq:bent_integral}
\end{equation}
while the path-length excess is
\begin{equation}
\Delta t_{2{\rm PN}}^{\rm (path)}
\;=\;\frac{1}{2c}\int\!\!\left(\frac{dx}{dz}\right)^{2}\!dz
\;+\;O(r_g^3),
\label{eq:path_integral}
\end{equation}
where $dx/dz$ is the $O(r_g)$ slope of the bent trajectory in
the straight-line coordinate~$z$.

\paragraph{The infra-red issue.}
Evaluating~\eqref{eq:bent_integral} or~\eqref{eq:path_integral}
over $(-\infty,+\infty)$ in either of the two conventions of
the preceding subsection produces \emph{divergent}
intermediate expressions (log or linear in the cutoff)---the
individual pieces are not separately finite, and the finite
physical answer emerges only after consistent finite-endpoint
regularization at emitter radius $r_1$ and receiver radius
$r_2$.  Finite-distance structure for the
\emph{refractive}-only master integral is noted in
Sec.~\ref{sec:finite} below; the analogous finite-distance
treatment for~\eqref{eq:bent_integral}--\eqref{eq:path_integral},
with matching between the near-lens perturbative region and
the asymptotic deflected-straight-line region, is the content
of the open calculation.

\paragraph{Consequence for the main theory.}
Until both pieces are closed in a single consistent
convention with matched asymptotics, neither the individual
coefficients of~\eqref{eq:bent2pn}--\eqref{eq:path2pn} nor the
total 2PN coefficient are asserted.  The main theory Shapiro section~\cite{Lotishvili2026}
correspondingly carries only the structural claim.

\section{Finite-distance endpoints: the refractive master integral}\label{sec:finite}
For emitter and receiver at finite radii $r_1,r_2\gg b$, the
straight-line limits of the refractive
integral~\eqref{eq:refr2pi} become $z\in[-z_1,z_2]$ with
$z_i=\sqrt{r_i^2-b^2}$.  The master integral is then
\begin{equation}
\int_{-z_1}^{z_2}\frac{dz}{b^2+z^2}
=\frac{1}{b}\left[\arctan\!\left(\frac{z}{b}\right)\right]_{-z_1}^{z_2}
=\frac{\pi}{b}-\frac{1}{r_1}-\frac{1}{r_2}+O\!\left(\frac{b^2}{r_i^3}\right),
\label{eq:finite}
\end{equation}
so the refractive 2PN coefficient is modified only by
corrections of relative order $O(b/r_i)$, which are a
theory-independent geometric factor.  For the Sgr~A$^\ast$
pulsar-timing geometry with $r_i\gtrsim 10^2 r_g$ and
$b\sim (10$--$30)\,r_g$ these are percent-level; for the solar
system they are negligible.  The analogous finite-endpoint
formulas for~\eqref{eq:bent_integral}--\eqref{eq:path_integral}
require the open calculation of Sec.~\ref{sec:open}.

\section{ISPG--GR 2PN difference (closed)}\label{sec:diff}

A finite, IR-stable statement is available about the
\emph{difference} between the ISPG and GR 2PN Shapiro
coefficients, even though each individual coefficient remains
open in the present decomposition (Sec.~\ref{sec:open}).

\subsection{Decomposition with shared 1PN trajectory}

Both theories produce the same $O(r_g)$ refractive index,
$n=1+2r_g/r+O(r_g^2)$, hence the same $O(r_g)$ light-bending
trajectory.  Writing the trajectory in flat-Cartesian
parameterization with $z$ along the asymptotic direction,
\begin{equation}
x(z)=b+r_g\,x_1(z)+r_g^2\,x_2(z)+O(r_g^3),
\label{eq:traj_exp}
\end{equation}
the 2PN expansion of the time-of-flight integral
\eqref{eq:dt_app6} reads
\begin{equation}
\Delta t_{2{\rm PN}}=\frac{r_g^2}{c}\int_{-L_1}^{+L_2}
\!\left[\,\frac{\alpha}{r_0(z)^2}
-\frac{2b\,x_1(z)}{r_0(z)^3}
+\tfrac12\,x_1'(z)^2\,\right]dz,
\label{eq:2pn_full}
\end{equation}
with $r_0(z)=\sqrt{b^2+z^2}$, $\alpha$ the coefficient of
$(r_g/r)^2$ in the refractive index, and $L_{1,2}$ the
asymptotic-direction distances of emitter and receiver from
the lens.  For the two metrics,
\begin{equation}
\alpha_{\rm ISPG}=2 \quad (\text{from } e^{2r_g/r}),
\qquad
\alpha_{\rm GR}=\tfrac{7}{4} \quad (\text{from }(1+r_g/(2r))^3/(1-r_g/(2r))).
\label{eq:alphas_app6}
\end{equation}
The remaining two integrands in~\eqref{eq:2pn_full} depend
only on the shared 1PN trajectory~$x_1$ and are therefore
\emph{identical} in ISPG and GR.

\subsection{Cancellation in the difference}

Subtracting the GR result from the ISPG result, the
$x_1$-dependent integrals cancel exactly, leaving
\begin{equation}
\Delta t_{2{\rm PN}}^{\rm ISPG}-\Delta t_{2{\rm PN}}^{\rm GR}
=\frac{r_g^2}{c}\,(\alpha_{\rm ISPG}-\alpha_{\rm GR})
 \int_{-L_1}^{+L_2}\!\frac{dz}{b^2+z^2}
=\frac{r_g^2}{4cb}\Bigl[\arctan\!\bigl(\tfrac{z}{b}\bigr)\Bigr]_{-L_1}^{+L_2},
\label{eq:db_finite}
\end{equation}
which evaluates in the asymptotic limit $L_1,L_2\to\infty$ to
\begin{equation}
\boxed{\;
\Delta t_{2{\rm PN}}^{\rm ISPG}-\Delta t_{2{\rm PN}}^{\rm GR}
=\frac{r_g^2}{cb}\cdot\frac{\pi}{4}\;}
\label{eq:db_inf}
\end{equation}
in the impact-parameter convention used uniformly
in~\eqref{eq:2pn_full} for both metrics.  The result is finite,
infra-red insensitive (the divergent path-length and
bent-ray pieces cancel by construction), and unaffected by
the convention dilemma of Sec.~\ref{sec:open} because any
$O(r_g)$ shift of $b$ acts identically on both theories and
drops out of the difference.

\subsection{Numerical verification}

A boundary-value formulation of the full Fermat ODE---shooting
for the bent trajectory between symmetric Cartesian
endpoints at $(z,x)=(\pm L,b)$ in each metric in turn, then
extracting the $O(r_g^2)$ coefficient by polynomial regression
in $r_g$---confirms the finite-$L$
formula~\eqref{eq:db_finite}.  Writing $\Delta B$ for the
dimensionless coefficient of $r_g^2/(cb)$ in the ISPG--GR
difference, with $b=1$:
\begin{center}
\begin{tabular}{rccc}
\toprule
$L/b$ & $\Delta B$ from~\eqref{eq:db_finite} & $\Delta B$ numerical & rel.\ error \\
\midrule
30 & 0.7687 & 0.7680 & 0.1\% \\
50 & 0.7754 & 0.7728 & 0.3\% \\
\bottomrule
\end{tabular}
\end{center}
At larger $L$ the integrated ODE error in the BVP shooting
accumulates and the numerics drift from the analytic curve;
the controlling statement is the closed
formula~\eqref{eq:db_finite} together with its
$L\to\infty$ limit~\eqref{eq:db_inf}.  The driver is
\texttt{MAIN/Python/Shapiro/check\_shapiro\_bvp.py}.

\subsection{Significance for the main theory}

The individual 2PN coefficients $B_{\rm ISPG}$ and $B_{\rm GR}$
are each IR-divergent at intermediate stages of the
decomposition~\eqref{eq:2pn_full} (Sec.~\ref{sec:open});
their absolute values require the open finite-endpoint
calculation flagged there.  The
\emph{difference}~\eqref{eq:db_inf}, however, is the sole
$O(r_g^2/(cb))$ Shapiro-delay correction that distinguishes
the bi-conformal ISPG metric from the GR Schwarzschild
geometry, and it is fixed without further model input.
This is the operationally relevant ISPG 2PN signature.

\section*{Takeaway}
The bi-conformal ISPG geometry produces an exact
\emph{refractive} 2PN Shapiro contribution
\[
\Delta t_{2{\rm PN}}^{\rm (refr)}=\frac{r_g^2}{cb}\,(2\pi),
\]
which is convention-independent and numerically verified.
The bent-ray and path-length contributions to the
\emph{individual} 2PN coefficient remain open
(Sec.~\ref{sec:open}), pending consistent finite-endpoint
regularization.  Their dependence on the shared $O(r_g)$
trajectory makes them identical in ISPG and GR, and the
corresponding ISPG--GR 2PN \emph{difference}
\[
\Delta t_{2{\rm PN}}^{\rm ISPG}-\Delta t_{2{\rm PN}}^{\rm GR}
=\frac{r_g^2}{cb}\cdot\frac{\pi}{4}
\]
is closed analytically and verified numerically
(Sec.~\ref{sec:diff}).  This is the operationally meaningful
ISPG signature at 2PN.

\end{document}
